% Pursuing Stacks (cotes 134-2 to 134-8): archive cote 134-6, site page 116 (PDF page 117), typescript page 434.
% typescript. Draft compared: ps5.tex:7561-7645.
% AI-transcribed by Claude Opus 5.5 (claude-opus-5-5, adaptive thinking, effort high), Message Batches API, batch ps3 (msgbatch_01XdFmuVVBAcgc7EFxdtCqsj), 2026-11-12; not yet checked by eye.
\scanpage{134-6/116}
\mspage{434}

\xout{(compare p. 403 formula (**} (this is formula (**) of page 403 \xout{for} with
P, \hand{$\mathcal{N}$} replaced by Q, \hand{$\mathcal{M}$}). It should be noted that the assumptions \inserted{on $\mathcal{N}$} made
in prop.4 are the weakest possible for the conclusions to hold\deleted{\xout{ ( as these}}
\deleted{(}for any k-additive small category P\deleted{)}, as these conclusions, in case
P = final category, just reduce to the assumptions\hand{.} \deleted{made}.

Here is the \inserted[typed]{outline} \xout{principle} of proof of prop.4. Using only stability
of \hand{$\mathcal{N}$} under small direct limits (besides k-additivity) we define a
canonical \xout{pairing} k-biadditive pairing
\[ (126) \qquad P^{\&} \mathrm{x} \hand{\mathcal{N}} \hand{\longrightarrow} P^{\&}_{\hand{\mathcal{N}}} , \qquad (L,x) \hand{\longmapsto} L\otimes_k x \overset{\mathrm{def}}{=} (a \hand{\mapsto} L(a)\otimes_k x ). \]
( NB I recall that $P^{\&}$ is defined as
\[ P^{\&} = \underline{Hom}_{\underline{Z}}(P^{o},Ab) \hand{\overset{\approx}{\longleftarrow}} \underline{Hom}_k(P^{o}, Ab_k ) , \]
here we interprete an object of $P^{\&}$ as a k-additive functor
\[ L : P^{o} \hand{\longrightarrow} Ab_k \; (\overset{\mathrm{def}}{=} (k\text{-}Mod)) \quad .) \]
The relevant fact here for objects of $P^{\&}_{\hand{\mathcal{N}}}$ of the type $L\otimes_k x$ is
\[ (127) \qquad Hom_{P^{\hand{\&}}_{\hand{\mathcal{N}}}}(L\otimes_k x,F) \hand{\simeq} Hom_{\hand{\mathcal{N}}}(x,Hom_k(L,F)) \hand{\simeq} Hom_{P^{\hand{\&}}}(L, Hom_{\deleted{\hand{\unsure{\mathcal{N}}}}}(x,F)) , \]
\margin{$P^{\&}_{\mathcal{N}}$}
where in the second term,
\[ Hom_k(L,F) \hand{\in} Ob \hand{\mathcal{N}} \]
is defined in a way dual to $F *_k L'$ (cf. comments above), using the
equivalence (dual to (*) above)
\[ (*') \qquad P^{\&}_{\hand{\mathcal{N}}} \hand{\simeq} \underline{Hom}_k{}^{\hand{!}}(P^{\& o}, \hand{\mathcal{N}}) , \]
which is defined only, however, when \hand{$\mathcal{N}$} is stable under small inverse
limits (hence the first isomorphism in (127) makes a sense ohly under
this extra assumption); \over{O}{o}n the other hand, in the third term in (127),
\[ Hom(x,F) \overset{\mathrm{def}}{=} (a \hand{\longmapsto} Hom_{\hand{\mathcal{N}}}(x,F(a)) \quad \text{in } P^{\hand{\&}} , \]
and the isomorphism between the first and third term in (127) makes a
sense and is defined without any extra assumption on \hand{$\mathcal{N}$}.

We leave to the reader to check (127) (where one is readily redu-%
ced to the case when L is an object a in P, using the \xout{exactness proper-}
commutation of the three functors obtained $P^{\& o} \hand{\longmapsto} Ab_k$ with small
inverse limits). \xout{When} It follows, when \hand{$\mathcal{N}$} is abelian:
\[ (128) \qquad L \text{ projective in } P^{\&} , \ x \text{ projective in } \hand{\mathcal{N}} \hand{\Rightarrow} L\otimes_k x \text{ proj. in } P^{\&}_{\hand{\mathcal{N}}} . \]

Assume now that any object of \hand{$\mathcal{N}$} is quotient of a projective one,
and let F be any object in $P^{\&}_{\hand{\mathcal{N}}}$. Formula (127) for L=a in P reduces

\margin{116}
